\section{Naiver Bayes}
\setcounter{aufgabe}{0}

% ============================================================
% AUFGABE 1: Satz von Bayes + totale Wahrscheinlichkeit
% ============================================================
\aufgabe{
Ein Betrieb sortiert \textbf{Lachs} und \textbf{Barsch}. In der Produktion sind
$70\,\%$ der Fische Lachse und $30\,\%$ Barsche. Ein Helligkeitssensor meldet
\textbf{hell}: Lachse sind zu $80\,\%$ hell, Barsche nur zu $20\,\%$.

\begin{enumerate}[label=\alph*)]
  \item Berechnen Sie über den Satz der totalen Wahrscheinlichkeit die
        Wahrscheinlichkeit $P(\text{hell})$.
  \item Ein Fisch wird als \textbf{hell} gemessen. Bestimmen Sie mit dem Satz
        von Bayes $P(\text{Lachs} \mid \text{hell})$.
  \item Wie wahrscheinlich ist es, dass dieser Fisch ein \textbf{Barsch} ist?
\end{enumerate}
}

\loesung{
\textbf{Gegebene Größen}
\[
\begin{aligned}
P(\text{Lachs}) &= 0{,}70 & P(\text{Barsch}) &= 0{,}30\\
P(\text{hell} \mid \text{Lachs}) &= 0{,}80 &
P(\text{hell} \mid \text{Barsch}) &= 0{,}20
\end{aligned}
\]

\begin{enumerate}[label=\alph*)]
  \item Totale Wahrscheinlichkeit für \textbf{hell}:
  \[
    P(\text{hell}) = P(\text{hell} \mid \text{Lachs})\,P(\text{Lachs})
                   + P(\text{hell} \mid \text{Barsch})\,P(\text{Barsch})
  \]
  \[
    P(\text{hell}) = 0{,}80 \cdot 0{,}70 + 0{,}20 \cdot 0{,}30
                   = 0{,}56 + 0{,}06 = 0{,}62.
  \]

  \item Satz von Bayes:
  \[
    P(\text{Lachs} \mid \text{hell})
      = \frac{P(\text{hell} \mid \text{Lachs})\,P(\text{Lachs})}{P(\text{hell})}
      = \frac{0{,}56}{0{,}62} \approx 0{,}903.
  \]

  \item Gegenwahrscheinlichkeit:
  \[
    P(\text{Barsch} \mid \text{hell}) = 1 - P(\text{Lachs} \mid \text{hell})
      = \frac{0{,}06}{0{,}62} \approx 0{,}097.
  \]
\end{enumerate}
}


% ============================================================
% AUFGABE 2: Naive Bayes mit diskreten Merkmalen
% ============================================================
\aufgabe{
Gegeben sei folgender Trainingsdatensatz mit den binären Merkmalen
\textbf{Länge} und \textbf{Helligkeit} sowie der Klasse (Lachs/Barsch).
\begin{center}
{\scriptsize
\renewcommand{\arraystretch}{1.25}
\setlength{\tabcolsep}{6pt}
\rowcolors{2}{SecondaryColor!12}{white}
\begin{tabular}{c c c c}
\rowcolor{PrimaryColor}
\textcolor{white}{\textbf{Nr.}} &
\textcolor{white}{\textbf{Länge}} &
\textcolor{white}{\textbf{Helligkeit}} &
\textcolor{white}{\textbf{Klasse}}\\
1 & lang & hell   & Lachs\\
2 & lang & hell   & Lachs\\
3 & lang & dunkel & Lachs\\
4 & lang & hell   & Lachs\\
5 & kurz & hell   & Lachs\\
6 & lang & hell   & Lachs\\
7 & kurz & dunkel & Barsch\\
8 & kurz & dunkel & Barsch\\
9 & lang & dunkel & Barsch\\
10& kurz & hell   & Barsch\\
\end{tabular}}
\end{center}

Ein neuer Fisch $B$ hat \textbf{Länge $=$ lang} und \textbf{Helligkeit $=$ hell}.
Klassifizieren Sie $B$ mit dem Naive-Bayes-Klassifikator.
\begin{enumerate}[label=\alph*)]
  \item Schätzen Sie die A-priori-Wahrscheinlichkeiten und die bedingten
        Wahrscheinlichkeiten aus den Häufigkeiten.
  \item Berechnen Sie für beide Klassen den Term
        $P(K)\cdot P(\text{Länge} \mid K)\cdot P(\text{Helligkeit} \mid K)$
        und bestimmen Sie die Klasse von $B$.
\end{enumerate}
}

\loesung{
Für den Vergleich genügt der Zähler $P(K)\cdot\prod_i P(M_i \mid K)$, da der
Nenner $P(B)$ für beide Klassen gleich ist.

\begin{enumerate}[label=\alph*)]
  \item Von $10$ Fischen sind $6$ Lachse und $4$ Barsche:
  \[
    P(\text{Lachs}) = \tfrac{6}{10} = 0{,}6,
    \qquad
    P(\text{Barsch}) = \tfrac{4}{10} = 0{,}4.
  \]
  Bedingte Wahrscheinlichkeiten:
  \[
    P(\text{lang} \mid \text{Lachs}) = \tfrac{5}{6},
    \qquad
    P(\text{hell} \mid \text{Lachs}) = \tfrac{5}{6},
  \]
  \[
    P(\text{lang} \mid \text{Barsch}) = \tfrac{1}{4},
    \qquad
    P(\text{hell} \mid \text{Barsch}) = \tfrac{1}{4}.
  \]

  \item Zähler je Klasse:
  \[
    P(\text{Lachs})\cdot P(\text{lang}\mid\text{Lachs})\cdot P(\text{hell}\mid\text{Lachs})
      = 0{,}6 \cdot \tfrac{5}{6} \cdot \tfrac{5}{6}
      = 0{,}6 \cdot \tfrac{25}{36} \approx 0{,}4167.
  \]
  \[
    P(\text{Barsch})\cdot P(\text{lang}\mid\text{Barsch})\cdot P(\text{hell}\mid\text{Barsch})
      = 0{,}4 \cdot \tfrac{1}{4} \cdot \tfrac{1}{4}
      = 0{,}4 \cdot \tfrac{1}{16} = 0{,}025.
  \]
  Wegen $0{,}4167 > 0{,}025$ wird $B$ als \textbf{Lachs} klassifiziert.

  Normiert ergibt sich:
  \[
    P(\text{Lachs} \mid B) = \frac{0{,}4167}{0{,}4167 + 0{,}025}
      \approx 0{,}943,
    \qquad
    P(\text{Barsch} \mid B) \approx 0{,}057.
  \]
\end{enumerate}
}


% ============================================================
% AUFGABE 3: Naive Bayes mit Gauß-Dichte
% ============================================================
\aufgabe{
Für die Merkmale \textbf{Helligkeit} und \textbf{Länge} wurden pro Klasse
Normalverteilungen mit Mittelwert $\mu$ und Standardabweichung $\sigma$
geschätzt:
\begin{center}
{\scriptsize
\renewcommand{\arraystretch}{1.25}
\setlength{\tabcolsep}{6pt}
\rowcolors{2}{SecondaryColor!12}{white}
\begin{tabular}{c c c c c}
\rowcolor{PrimaryColor}
\textcolor{white}{\textbf{Klasse}} &
\textcolor{white}{\textbf{$\mu_{\text{Hell}}$}} &
\textcolor{white}{\textbf{$\sigma_{\text{Hell}}$}} &
\textcolor{white}{\textbf{$\mu_{\text{Länge}}$}} &
\textcolor{white}{\textbf{$\sigma_{\text{Länge}}$}}\\
Lachs  & 6 & 2 & 12 & 3\\
Barsch & 9 & 2 & 8  & 3\\
\end{tabular}}
\end{center}

Die Klassen sind gleich häufig, also $P(\text{Lachs}) = P(\text{Barsch}) = 0{,}5$.
Ein neuer Fisch $B$ hat \textbf{Helligkeit $= 7$} und \textbf{Länge $= 10$}.
Klassifizieren Sie $B$ mit dem Naive-Bayes-Klassifikator über die
Gauß-Dichtefunktion
\[
  f(x) = \frac{1}{\sigma\sqrt{2\pi}}\,
         \exp\!\Bigl(-\frac{(x-\mu)^2}{2\sigma^2}\Bigr).
\]
}

\loesung{
Wegen der Unabhängigkeitsannahme gilt
$P(K)\cdot f(\text{Hell}{=}7 \mid K)\cdot f(\text{Länge}{=}10 \mid K)$.

\medskip
\textbf{Dichtewerte Lachs}
\[
  f(\text{Hell}{=}7 \mid \text{Lachs})
    = \frac{1}{2\sqrt{2\pi}}\,\exp\!\Bigl(-\frac{(7-6)^2}{2\cdot 2^2}\Bigr)
    = 0{,}1995 \cdot e^{-0{,}125} \approx 0{,}1760,
\]
\[
  f(\text{Länge}{=}10 \mid \text{Lachs})
    = \frac{1}{3\sqrt{2\pi}}\,\exp\!\Bigl(-\frac{(10-12)^2}{2\cdot 3^2}\Bigr)
    = 0{,}1330 \cdot e^{-0{,}2222} \approx 0{,}1065.
\]

\textbf{Dichtewerte Barsch}
\[
  f(\text{Hell}{=}7 \mid \text{Barsch})
    = \frac{1}{2\sqrt{2\pi}}\,\exp\!\Bigl(-\frac{(7-9)^2}{2\cdot 2^2}\Bigr)
    = 0{,}1995 \cdot e^{-0{,}5} \approx 0{,}1210,
\]
\[
  f(\text{Länge}{=}10 \mid \text{Barsch})
    = \frac{1}{3\sqrt{2\pi}}\,\exp\!\Bigl(-\frac{(10-8)^2}{2\cdot 3^2}\Bigr)
    = 0{,}1330 \cdot e^{-0{,}2222} \approx 0{,}1065.
\]

\medskip
\textbf{Zähler je Klasse}
\[
  P(\text{Lachs})\cdot 0{,}1760 \cdot 0{,}1065
    = 0{,}5 \cdot 0{,}01874 \approx 0{,}00937,
\]
\[
  P(\text{Barsch})\cdot 0{,}1210 \cdot 0{,}1065
    = 0{,}5 \cdot 0{,}01288 \approx 0{,}00644.
\]
Wegen $0{,}00937 > 0{,}00644$ wird $B$ als \textbf{Lachs} klassifiziert.

Normiert:
\[
  P(\text{Lachs} \mid B) \approx \frac{0{,}00937}{0{,}00937 + 0{,}00644}
    \approx 0{,}593,
  \qquad
  P(\text{Barsch} \mid B) \approx 0{,}407.
\]
}